% 数学ノート集・共通テンプレート v2（積分論2026年度第1部の版面を参考）
% A5判。章内で定義・定理・例・注意等の番号を共通化。
% コンパイル： platex -interaction=nonstopmode math_note_template_v2.tex (2回)
%              dvipdfmx math_note_template_v2.dvi
\documentclass[a5paper,10pt,onecolumn]{jsarticle}
\usepackage[dvipdfmx,top=22mm,bottom=20mm,inner=19mm,outer=18mm,headheight=13pt,headsep=6mm]{geometry}
\usepackage{amsmath,amssymb,amsthm,mathtools}
\usepackage{enumitem,fancyhdr,titlesec}
\usepackage[dvipdfmx,bookmarks=true,colorlinks=false]{hyperref}
\usepackage{pxjahyper}

% ----- 差し替える文書情報 -----
\newcommand{\NoteTitle}{測地線と指数写像}
\newcommand{\NoteSeries}{数学ノート集}

% ----- 文字・組版 -----
% 参考PDFの質素な明朝＋Computer Modernの印象を引き継ぎながら，
% A5判で窮屈にならないよう本文の行間を少し確保する．
\linespread{1.13}
\setlength{\parindent}{1zw}
\setlength{\parskip}{0pt}
\setlength{\abovedisplayskip}{6pt plus 2pt minus 2pt}
\setlength{\belowdisplayskip}{6pt plus 2pt minus 2pt}
\setlength{\abovedisplayshortskip}{3pt plus 2pt minus 1pt}
\setlength{\belowdisplayshortskip}{4pt plus 2pt minus 1pt}
\setlist[enumerate]{topsep=3pt,itemsep=1pt,parsep=0pt,leftmargin=2.4zw}
\titleformat{\section}{\normalfont\normalsize\bfseries}{\thesection}{0.8zw}{}
\titlespacing*{\section}{0pt}{2.0ex plus .4ex}{.9ex}
\titleformat{\subsection}{\normalfont\normalsize\bfseries}{\thesubsection}{0.7zw}{}
\titlespacing*{\subsection}{0pt}{1.5ex plus .3ex}{.6ex}
\raggedbottom

% ----- 柱とページ番号：ページ上部、罫線なし -----
\fancypagestyle{mathnote}{%
  \fancyhf{}%
  \fancyhead[L]{\footnotesize \NoteSeries\ （\NoteTitle）}%
  \fancyhead[R]{\footnotesize \thepage}%
  \renewcommand{\headrulewidth}{0pt}%
  \renewcommand{\footrulewidth}{0pt}%
}
\pagestyle{mathnote}

% ----- 節内通し番号の定義・定理群 -----
\newtheoremstyle{mathnotetheorem}
 {4.5pt}{4.5pt}{\normalfont}{}%
 {\bfseries}{}{0.65zw}%
 {\thmname{#1}\thmnumber{ #2}\thmnote{（#3）}}
\theoremstyle{mathnotetheorem}
\newtheorem{theorem}{定理}[section]
\newtheorem{definition}[theorem]{定義}
\newtheorem{lemma}[theorem]{補題}
\newtheorem{proposition}[theorem]{命題}
\newtheorem{corollary}[theorem]{系}
\newtheorem{example}[theorem]{例}
\newtheorem{remark}[theorem]{注意}
\renewcommand{\proofname}{証明}
\renewcommand{\qedsymbol}{$\square$}
\makeatletter
\renewenvironment{proof}[1][\proofname]{%
  \par\pushQED{\qed}\normalfont\topsep4\p@\@plus2\p@\relax
  \trivlist\item[\hskip\labelsep\bfseries #1]\ignorespaces
}{\popQED\endtrivlist\@endpefalse}
\makeatother

% ----- 使うノートで適宜拡張する記号 -----
\newcommand{\R}{\mathbb{R}}
\newcommand{\norm}[1]{\lVert #1\rVert}
\newcommand{\abs}[1]{\lvert #1\rvert}

\begin{document}
\section*{読了後にできること}
本稿を読むための前提は，Levi--Civita接続，測地線方程式，指数写像の定義．
本稿の具体的な到達課題は，定曲率のモデルで測地線と指数写像を実際に求め，その定義域を調べる．


\section{測地線の局所理論}
\begin{definition}
$(M,g)$ のLevi-Civita接続を $\nabla$ とする．
滑らかな曲線 $\gamma$ が
$\nabla_{\dot\gamma}\dot\gamma=0$ を満たすとき
測地線という．
\end{definition}
\begin{proposition}[測地線方程式]\label{prop:geod}
局所座標 $x^k(t)$ において
\[
 \ddot x^k+\sum_{i,j}\Gamma^k_{ij}(x)\dot x^i\dot x^j=0
 \quad(k=1,\dots,n)
\]
である．任意の初期位置 $p$ と初期速度 $v$ に対して
短時間の測地線がただ一つ存在する．
\end{proposition}
\begin{proof}
$\nabla_{\dot\gamma}\dot\gamma$ の座標成分を
接続のLeibniz則で展開すると表示式を得る．
これは $(x,\dot x)$ について滑らかな一階連立常微分方程式なので，
Picard--Lindel\"ofの定理と滑らかな依存性から
短時間の解の存在一意性を得る．
\end{proof}
\begin{proposition}[測地線の速さは一定]
測地線 $\gamma$ では $|\dot\gamma|_g$ は一定である．
\end{proposition}
\begin{proof}
計量保存性から
\[
 \frac d{dt}g(\dot\gamma,\dot\gamma)
 =2g(\nabla_{\dot\gamma}\dot\gamma,\dot\gamma)=0
\]
となる．
\end{proof}
\section{指数写像}
\begin{definition}
$\gamma_v(0)=p$，$\dot\gamma_v(0)=v$ を満たす測地線が
時刻1まで定義される $v\in T_pM$ の集合を $\mathcal D_p$ とし，
\[
 \exp_p(v)=\gamma_v(1)
\]
を指数写像という．
\end{definition}
\begin{lemma}[時間の再尺度化]
$\gamma_{cv}(t)=\gamma_v(ct)$ が両辺の定義される範囲で成立する．
従って $\exp_p(tv)=\gamma_v(t)$ である．
\end{lemma}
\begin{proof}
$\eta(t)=\gamma_v(ct)$ に対して
$\dot\eta(0)=cv$，
$\nabla_{\dot\eta}\dot\eta=c^2
\nabla_{\dot\gamma_v}\dot\gamma_v=0$．
測地線の一意性から $\eta=\gamma_{cv}$ である．
\end{proof}
\begin{theorem}[正規座標の存在]\label{thm:normal}
$0\in T_pM$ のある開近傍 $V$ で
$\exp_p:V\to M$ は $p$ の開近傍への微分同相となる．
\end{theorem}
\begin{proof}
測地線ODEの滑らかな依存性から
$\exp_p$ は $0$ 近傍で滑らかである．
時間再尺度化により
$\exp_p(tv)=\gamma_v(t)$ だから
\[
 (d\exp_p)_0v=
 \left.\frac d{dt}\right|_0\exp_p(tv)
 =\dot\gamma_v(0)=v.
\]
従って $(d\exp_p)_0$ は恒等写像であり，
逆写像定理から結論が従う．
\end{proof}
\section{Gaussの補題}
\begin{theorem}[Gaussの補題]\label{thm:gauss}
$v\in T_pM$ を指数写像の定義域の点とし，
$w\in T_pM$ とする．このとき
\[
 g_{\exp_p(v)}
 \bigl((d\exp_p)_v v,(d\exp_p)_v w\bigr)
 =g_p(v,w)
\]
が成り立つ．
\end{theorem}
\begin{proof}
$F(s,t)=\exp_p(t(v+sw))$ を考える．
$t$ 方向の各曲線はアフィンパラメータ付き測地線である．
$\partial_sF=V$，$\partial_tF=T$ と置くと
$\nabla_TT=0$ かつ捩率零から
$\nabla_TV=\nabla_VT$ である．
従って
\[
 \partial_t g(T,V)
 =g(T,\nabla_TV)
 =g(T,\nabla_VT)
 =\tfrac12\partial_sg(T,T).
\]
測地線の速さは一定であり
$g(T,T)=g_p(v+sw,v+sw)$ なので，
$s=0$ で右辺は $g_p(v,w)$ となる．
$t=0$ では $V(0,0)=0$ なので
$g(T,V)=0$，従って
$g(T,V)=t\,g_p(v,w)$．
$t=1$ で $T=(d\exp_p)_v v$，
$V=(d\exp_p)_v w$ となるから結論を得る．
\end{proof}
\section{局所最短性}
\begin{theorem}[正規球内の最短性]\label{thm:min}
$p$ の十分小さな正規球 $U=\exp_p(B_r(0))$ に対し，
任意の $q=\exp_p(v)\in U$ を結ぶ放射測地線
$t\mapsto\exp_p(tv)$ は，
$U$ に含まれる区分的滑らかな曲線の中で長さが最小である．
その長さは $|v|_g$ となる．
\end{theorem}
\begin{proof}
$U\setminus\{p\}$ の曲線
$\sigma(t)=\exp_p(\rho(t)u(t))$ を考える．
$|u(t)|_p=1$ とする．
Gaussの補題から
$(d\exp_p)_{\rho u}u$ と
$(d\exp_p)_{\rho u}u'$ は直交するので
\[
 |\dot\sigma|_g^2
 =|\rho'|^2+
 |(d\exp_p)_{\rho u}(\rho u')|_g^2
 \ge|\rho'|^2.
\]
従って
$L(\sigma)\ge\int|\rho'|dt\ge
 |\rho(b)-\rho(a)|$．
始点が $p$ である一般の曲線は
$p$ を離れた部分へこの評価を適用し
極限を取ればよい．
終点が $\exp_p(v)$ なら
$L(\sigma)\ge|v|_g$．
放射測地線は一定速度 $|v|$ なので長さがちょうど
$|v|$ であり最短となる．
\end{proof}
\begin{remark}
上の定理は曲線が正規球内にある場合である．
大域的な距離の最小化には，近傍をさらに縮めた
強凸正規近傍を使用する．
\end{remark}

\section{到達課題の完全解答}
\noindent\textbf{球面上の指数写像の完全な計算．}
\begin{proposition}[単位球面の測地線]
$S^n\subset\mathbb R^{n+1}$ に標準計量を入れる．
$p\in S^n$，$v\in T_pS^n$ とすると，
$v\ne0$ に対し
\[
\gamma_v(t)=\cos(|v|t)p+
\frac{\sin(|v|t)}{|v|}v
\]
は初期条件 $\gamma_v(0)=p$，
$\dot\gamma_v(0)=v$ を満たす測地線である．
よって
\[
\exp_p(v)=\cos|v|\,p+
\frac{\sin|v|}{|v|}v
\]
となり，$v=0$ では $\exp_p(0)=p$ とする．
\end{proposition}
\begin{proof}
$p\cdot v=0$，$|p|=1$ だから
$|\gamma_v(t)|=1$．
また
$\gamma_v''(t)=-|v|^2\gamma_v(t)$ は
球面の法線方向であるため
接成分である
$\nabla_{\dot\gamma_v}\dot\gamma_v$ は零となる．
初期条件と測地線の一意性から結論が従う．
\end{proof}
\begin{corollary}[球面の指数写像の最初の特異点]
$n\ge2$ とする．任意の単位ベクトル $u\in T_pS^n$ に対し，
$\exp_p(\pi u)=-p$ であり，
$(d\exp_p)_{\pi u}$ は単射でない．
一方，$0\le t<\pi$ では $(d\exp_p)_{tu}$ は同型である．
従って各放射半直線上で最初に微分が特異となるのは長さ $\pi$ の点である．
\end{corollary}
\begin{proof}
$v=ru$（$r>0$, $|u|=1$）とすると，
$\exp_p(v)=\cos r\,p+\sin r\,u$ である．
$w\in T_pS^n$ を $u$ に直交するベクトルとする．
このとき $r$ の $w$ 方向の微分は零なので
\[
 (d\exp_p)_{ru}(w)=\frac{\sin r}{r}\,w,
 \qquad
 (d\exp_p)_{ru}(u)=-\sin r\,p+\cos r\,u.
\]
後者は長さ $1$ であり，前者と直交する．
$0<r<\pi$ なら $\sin r/r\ne0$ だから
$T_pS^n=\mathbb Ru\oplus u^\perp$ 上で微分は単射であり，
両辺の次元が $n$ なので同型である．
$r=\pi$ では $\sin r=0$ だから，
$n\ge2$ より存在する非零の $w\in u^\perp$ が微分の核に入る．
また $\exp_p(\pi u)=-p$ は表示式から従う．
$r=0$ では定理\ref{thm:normal}により
$(d\exp_p)_0$ は同型である．
\end{proof}
\begin{remark}
$n=1$ の場合は指数写像 $T_pS^1\cong\mathbb R\to S^1$ が
各点で局所微分同相であり，特異点を持たない．
ただし $\exp_p(\pi u)=\exp_p(-\pi u)=-p$ であり，
大域的には単射ではない．
\end{remark}
\noindent\textbf{到達確認．}
$S^2$ の北極 $p=(0,0,1)$ と
$v=(a,0,0)$ に対し $\exp_p(v)$ を求めよ．
\par\noindent\textbf{解答．}
$a\ne0$ なら
$\exp_p(v)=(\sin|a|\,a/|a|,0,\cos|a|)$
であり，$a=0$ なら北極となる．
\end{document}
